\chapter{The Texture Problem}

\section{A Wall That Changes Every Time You Look At It}

Players of AI-generated games have a name for a particular kind of small, persistent disappointment. You turn a corner, glance at a stone wall, keep moving, and a moment later, for whatever reason, look back. The wall is still there, still recognizably a wall, still built from the same kind of stone. But the moss has moved. The crack that ran diagonally through the third course of blocks is gone, or has migrated somewhere else entirely, or has been joined by a second crack that was never there a moment ago. Nothing about this is a rendering glitch in the ordinary sense; every individual frame looks fine. The problem only shows up in the difference between frames, in the fact that the system generating the wall never actually committed to a specific wall in the first place. It committed to a style of wall, and produced a fresh, independent, individually convincing instance of that style each time it was asked.

This is not a hypothetical concern raised for the sake of argument. It is a routine, observed behavior of the generative systems currently used to build interactive AI worlds, and it is worth taking seriously exactly because each individual output is so rarely at fault. A single rendering of the wall, examined on its own, usually looks correct, coherent, and plausible. The failure is invisible at the level of any one frame and only becomes visible in the relationship between frames, which is precisely the level at which most evaluation of generative systems does not look.

\section{Rendering Versus Remembering}

It is worth being precise about what has actually gone wrong, because the natural first description, that the system produced two different walls, understates the problem. The system did not produce two different walls. It never produced a wall, in the sense of a specific, individuated, continuing thing, at all. What it produced, twice, was a rendering: a plausible visual instance of "stone wall," generated fresh from a learned distribution over what stone walls tend to look like, with no reference to anything that happened the first time this particular prompt, or this particular scene, was rendered before.

Rendering, in this sense, and remembering are different operations, and the difference is not a matter of degree. Rendering asks a model what a stone wall plausibly looks like, and the model answers from its training distribution, drawing a fresh sample each time it is asked. Remembering asks what this stone wall, the specific one that has already been established, looks like, and answers by consulting a record of what was previously committed rather than by sampling anew. A system built entirely from rendering can be extraordinarily good at its actual job, producing individually beautiful, physically plausible, stylistically consistent images, and still have no mechanism whatsoever for remembering, because nothing about producing a good sample from a distribution requires consulting what any previous sample from that distribution happened to be.

\section{A Problem Bigger Than Graphics}

It would be a mistake to file this under graphics or rendering quality, because the same failure shows up wherever a generative system is asked to maintain something across more than one act of generation, regardless of modality. A video model asked to continue a shot will often drift, subtly reinterpreting a character's face, a room's layout, or an object's color across a cut it was supposed to be maintaining continuity through, for exactly the reason the wall drifts: nothing forces frame two to be a continuation of frame one rather than an independent, merely plausible successor to it. A language model maintaining a long conversation will, past a certain length, quietly contradict a detail it introduced many turns earlier, not because it has access to contradicting information, but because the detail was never stored anywhere retrievable; it was generated once, used briefly, and then effectively forgotten the moment it left the portion of context the model was actively attending to. A game character built from a language model will cheerfully claim, on the tenth conversation, to be meeting the player for the first time, having "forgotten" nine previous conversations not through any failure of memory in the ordinary sense but because nothing about how the character speaks was ever connected to a persistent record of what it had previously said.

These are not three unrelated bugs occurring in three different systems. They are the same architectural gap, showing up in whatever modality happens to be under discussion: a system built to render plausible outputs, asked to also maintain a persistent identity across more than one output, and given no mechanism to do so beyond hoping that plausibility and consistency happen to coincide.

\section{Why More Scale Does Not Fix This}

It is tempting to treat this as a problem scale will eventually solve: a larger model, trained on more data, conditioned on a longer context window, will surely learn to be more consistent, the way larger models have learned to be more fluent, more knowledgeable, more capable along nearly every other axis that has been thrown at them. This chapter will not attempt to prove that scale cannot help; it plainly can, and a sufficiently long context window genuinely does reduce how often a model contradicts something it said recently enough to still be attending to. But reduction is not elimination, and the reason is architectural rather than a matter of degree. A context window, however long, is still a buffer of recent inputs the model conditions on, not a persistent record the model is obligated to remain faithful to. Nothing about attending to more tokens changes the underlying operation from rendering into remembering; it only enlarges the window within which rendering happens to look like remembering, because enough of the relevant detail is still nearby enough to influence the next plausible sample. Push the wall far enough outside that window, long enough ago in the conversation, far enough back in the game session, and the drift this chapter opened with returns, precisely because it was never actually solved, only outrun for as long as the buffer lasted. This book will return to this argument in more detail in the chapters immediately following this one; for now it is enough to notice that the problem this chapter has described survives every scale of system that has been built to render rather than to remember, and shows no sign, in principle, of a scale at which it would stop.

\section{What a Real Solution Would Need}

It is worth naming, in outline and without yet attempting to build it, what a system would actually need in order to avoid this failure, because the shape of the requirement turns out to matter more than any specific technique for satisfying it. Such a system would need something that continues to exist between acts of generation, independent of whether anything is currently observing it, so that there is something for a second observation to be checked against rather than merely a second independent sample. It would need to treat observation itself differently: not as an occasion to produce a fresh, plausible output, but as an occasion to read from whatever has persisted, producing new content only where genuinely nothing has been established yet. And it would need this discipline to hold not as a matter of the model's trained preferences, which can always be overridden by a sufficiently unusual prompt or a sufficiently long gap, but as a matter of the system's own architecture, in the same sense that a database's consistency does not depend on the application querying it happening to ask politely.

This book is an extended answer to what a system built this way would actually look like: what it would mean for something to persist, how observation could be redefined as reading rather than inventing, and what would have to be true of a system's internals for this discipline to be architectural rather than aspirational. None of that begins in this chapter. This chapter's task was only to make the problem impossible to unsee, in enough different guises, that a reader arrives at the rest of the book already convinced that the problem is real, structural, and not obviously solved by anything already on the table.

\section{Looking Ahead}

The wall, the video cut, the forgetful character are all symptoms of a single underlying fact about how current systems are built: they have no persistent internal state at all, only a context they condition on and a policy for producing the next plausible output given that context. Chapter 2 examines this fact directly, asking what, if anything, current architectures do have in place of a persistent world, and why what they have, however sophisticated, still falls short of the thing this book is going to insist on calling a world.
